\section{Choosing the weights for one shift}\label{sec:weights}

At a fixed shift, a weight sends $(1,0,h)$ to a thick place when its
index belongs to a prescribed translate. We choose the indices so that
every translate contains nearly the same number.

\subsection{Counts on translates}

Fix $0<\varepsilon<1/52$ and use the field at $x$ from
Proposition~\ref{prop:field-window}. Fix the set $T_2\subset H$ of
thick indices before choosing the weight indices $T\subset H$.
Only the size bound $N_2=|T_2|\ge n x^{-\varepsilon}$ enters the count.
The construction in \citep[Proposition~5.17]{adriazola-grh} gives
such a set. Write
\[
  \delta_2=N_2/n,\qquad p=x^{-\varepsilon/8},\qquad
  N_T=|T|,\qquad C_s=|T\cap(T_2+s)|.
\]

\begin{lemma}\label{lem:weight-count}
For sufficiently large $x$ there is a set $T$ with
$pn/2\le N_T\le 2pn$ and
\begin{equation}\label{eq:weight-count}
   \max_{s\in H}\left|\frac{C_s}{N_T\delta_2}-1\right|\le4\kappa,
   \qquad \kappa=\left(\frac{12\log n}{pN_2}\right)^{1/2}.
\end{equation}
Given prime and orbit data as in \citep[Proposition~5.20]{adriazola-grh},
one choice of $T$ satisfies both \eqref{eq:weight-count} and the
weighted bound on exceptional Frobenius orbits in that proposition.
\end{lemma}

\begin{proof}
Include each $t\in H$ independently with probability $p$. For fixed
$s$, the count $C_s$ has binomial law with mean $\mu=pN_2$.
Since $\mu\ge x^{2+55\varepsilon/8}$, we have $\kappa\le1/2$ for
large $x$. The Chernoff bound gives
\[
  \Pr\!\left(\max_{s\in H}|C_s-\mu|>\kappa\mu\right)
  \le 2n\exp(-\kappa^2\mu/3)=2n^{-3}.
\]
Each $t\in T$ occurs in $N_2$ translates, so
$n^{-1}\sum_s C_s=N_T\delta_2$. On the complementary event both
$C_s$ and $N_T\delta_2$ lie in $[(1-\kappa)\mu,(1+\kappa)\mu]$.
This gives the size bound and \eqref{eq:weight-count}. For the trivial
target character the proof of the cited proposition uses the same
Bernoulli law. Its size and orbit bounds hold together with probability
at least $2/3-2/7>2n^{-3}$ for large $x$.
\end{proof}

\subsection{The excess at one shift}

Put $b=\exp(x^\varepsilon)$, $Z=\exp(x^{\varepsilon/3})$ and
$\omega=\exp(-x^{\varepsilon/8})$. For $u_0\in H$ call the places
$(0,1,s)$ with $s\in T_2+u_0$ thick and all other places in $\Phi$
thin. Let $m\ge1$ and write $w_j(\rho)>0$ for the nodes at this shift,
$1\le j\le m$. Assume
\begin{equation}\label{eq:weight-ranges}
  |a_t/b-x|\le x^{-1},\qquad
  \begin{cases}
    b\le w_j(\rho)\le b(1+x^{-2}) & \text{at thick places},\\
    w_j(\rho)\le 2xb/Z & \text{at thin places}.
  \end{cases}
\end{equation}
These inequalities follow from the numerical ranges in
\citep[Theorem~A(i) and Proposition~5.16]{adriazola-grh}.

\begin{proposition}\label{prop:one-shift}
Take $T$ as in Lemma~\ref{lem:weight-count} and put
$Q=N_Tmb\delta_2/2$. For any shift $u_0$ and weights and nodes
satisfying \eqref{eq:weight-ranges}, the excess of
\S\ref{sec:accounting} satisfies
\begin{equation}\label{eq:one-shift}
  \frac{D_{u_0}}{Q}=4\sum_{h\in H}F_x(1,0,h)+o(1).
\end{equation}
For fixed $\varepsilon$ and window exponent $J$, the error tends to
zero uniformly in these choices as $x\to\infty$.
\end{proposition}

\begin{proof}
Write $K_y(g)=F_y(g)/y$ and let $\mu_r(g)$ be the probability
in \eqref{eq:field-frobenius} for the prime power $r$. From the
definition of $F_y$,
\[
  K_y(g)=c_g-\frac{\log y}{4n}
     +\sum_{r\ \mathrm{prime\ power}}\Lambda(r)\mu_r(g)
          \left(\frac1r-\frac1y\right)_+ ,
\]
where $c_g$ collects the terms independent of $y$. Conjugate
characters pair in $F_y$, so $K_y$ is real. Put
$\|K_y\|_1=\sum_{g\in\mathcal G}|K_y(g)|$. Between consecutive
integers, $\sum_g\mu_r(g)=1$ and Chebyshev's bound give
\[
  \|K'_y\|_1\le\frac1y+\frac1{y^2}\sum_{r\le y}\Lambda(r)
    \ll \frac1y.
\]
Each summand is continuous at $y=r$. Integrating from $x$ to $Z$
and applying \eqref{eq:field-large-arguments} at $Z$ gives
\begin{equation}\label{eq:weight-variation}
  \|K_x\|_1\ll 4n\omega+\log(Z/x)\ll 1+\log Z.
\end{equation}

For $g\notin\Phi$ put $y_{jtg}=a_t/w_j(p_tg)$. The finite-place sum
in \S\ref{sec:accounting} gives the identity
\begin{equation}\label{eq:weight-excess}
  D_{u_0}=2\sum_{j,t}a_t\sum_{g\notin\Phi}K_{y_{jtg}}(g).
\end{equation}
For each $j,t$ the prime-power terms sum to the finite-place
contribution \citep[Proposition~3.3]{adriazola-grh}
\[
   S_{jt}=2\sum_{g\notin\Phi}\ \sum_{r\ \mathrm{prime\ power}}
        \Lambda(r)\mu_r(g)\bigl(a_t/r-w_j(p_tg)\bigr)_+ .
\]
The map $p_t$ sends the complement of $\Phi$ bijectively onto $\Phi$,
so the logarithms contribute $-a_t\log(a_t/\bar w_j)$ for each $j,t$.
Summing the characters over that complement gives $f$ for the trivial
character, $-f$ for $\chi_{-d}$ and zero for the others.
Thus $2\sum_{g\notin\Phi}c_g=1+\gamma-\lambda_1$, as in
\eqref{eq:account-main}.

At a thin place $y_{jtg}\ge Z/3$ and hence
$|K_{y_{jtg}}(g)|\le\omega$ by \eqref{eq:field-large-arguments}.
At a thick place $y_{jtg}$ lies in $I=[x-2/x,x+2/x]$ by
\eqref{eq:weight-ranges}. For fixed $j,t$ the derivative estimate gives
\[
 \sum_{p_tg\ \mathrm{thick}}
   |K_{y_{jtg}}(g)-K_x(g)|
 \le \int_I\|K'_v\|_1\,dv\ll x^{-2}.
\]
Discard the thin terms in \eqref{eq:weight-excess} and replace
$a_t K_{y_{jtg}}(g)$ at the thick places by $bxK_x(g)=bF_x(g)$.
After division by $Q$, the total error is
\[
  \ll \delta_2^{-1}\bigl(fx\omega+x^{-1}
                      +x^{-1}\|K_x\|_1\bigr)=o(1)
\]
by \eqref{eq:weight-variation} and $\delta_2\ge x^{-\varepsilon}$.

Now $p_t(1,i,h)=(0,i+1,h+t)$ is thick just when $i=0$ and
$t\in T_2+u_0-h$. The thick terms therefore give
\[
  \frac{D_{u_0}}{Q}
    =4\sum_{h\in H}\frac{C_{u_0-h}}{N_T\delta_2}F_x(1,0,h)+o(1).
\]
Replacing each ratio by $1$ costs at most $16\kappa x\|K_x\|_1$.
Since $pN_2\ge x^{2+55\varepsilon/8}$ and $\log Z=x^{\varepsilon/3}$,
\eqref{eq:weight-variation} gives $\kappa x\|K_x\|_1=o(1)$.
\end{proof}

Equations \eqref{eq:field-projection}, \eqref{eq:field-window-bound}
and \eqref{eq:one-shift} give
\[
  G(x)\le \frac{D_{u_0}}{Q}+3x^{1/2+\varepsilon}+o(1).
\]
To bound $D_{u_0}$ by Proposition~\ref{prop:accounting}, the arithmetic
construction must realize its seven inequalities for this $T$ and one
shift $u_0$ with the weights and nodes in \eqref{eq:weight-ranges}.
